\begin{proof}[Proof of \cref{obj:thm:two-cell-calibration}]
\begin{enumerate}
\item
By \cref{obj:lem:measurable-kernel-radon-nikodym}, the measurable-density hypothesis of \cref{obj:lem:dimension-free-private-two-point} is satisfied. The latter supplies a universal constant \(c_1>0\) such that, for every allowed two-cell resource tuple and sampling scheme, and every \(\mathcal C\in\{\mathrm{NI},\mathrm{SI}\}\),
\[
\mathfrak R_{\mathcal C}(n,2,\varepsilon)
\ge c_1\min\left\{1,\frac1{n\varepsilon^2}\right\},
\qquad
\mathfrak H_{\mathcal C}(n,2,\varepsilon)
\ge c_1\min\left\{1,\frac1{\sqrt n\,\varepsilon}\right\}.
\]
Set
\[
c=\min\{c_1,1\},
\qquad C_0=1000.
\]
Then \(0<c\le1\le C_0<\infty\).

Fix \(n\ge2\), \(0<\varepsilon\le1\), and a sampling scheme satisfying
\cref{obj:ass:iid-people,obj:ass:independent-randomness}. Define the effective private resource by
\[
t=n\varepsilon^2>0.
\]
All expectations and probabilities below use the decision law induced by this sampling scheme and the specified protocol. We first construct simultaneous risk and interval witnesses, splitting according to \(t\le1\) or \(t>1\).
% lean: CausalSmith.Stat.LdpOptvalueUniformFrontier.two_cell_calibration

\item
Suppose \(t\le1\). Since
\[
(\sqrt n\,\varepsilon)^2=t,
\qquad 0<\sqrt n\,\varepsilon\le1,
\]
the two rate factors satisfy
\[
\min\left\{1,\frac1{n\varepsilon^2}\right\}
=
\min\left\{1,\frac1{\sqrt n\,\varepsilon}\right\}
=1.
\]
Use a singleton public seed and let every participant send the same fixed message. Set
\[
T(\mathbf Z,R,U)=\frac12,
\qquad
I(\mathbf Z,R,U)=\left[\frac14,\frac34\right].
\]
The message law is constant in the original record and history, so it satisfies full-record privacy and noninteraction. The estimate and interval endpoints are constant Borel functions.

For every \(P\in\mathcal V_2\), the interior means in
\cref{obj:def:causal-model} give
\[
\frac14\le
\max\{\mu_{0j}(P),\mu_{1j}(P)\}
\le\frac34
\qquad(j\in[2]).
\]
Averaging these inequalities yields
\[
\frac14\le
V(P)=\frac12\sum_{j=1}^{2}
\max\{\mu_{0j}(P),\mu_{1j}(P)\}
\le\frac34.
\]
Consequently,
\[
\mathbb E[(T-V(P))^2]
=\left(\frac12-V(P)\right)^2
\le\frac1{16},
\qquad
\mathbb P\{V(P)\in I\}=1,
\qquad
\mathbb E[\operatorname{length}I]=\frac12.
\]
These bounds attain both rate factors with comparison constant \(1\). Enlarging that constant to \(1000\) preserves the same witnesses and coverage.
% lean: CausalSmith.Stat.LdpOptvalueUniformFrontier.two_cell_saturated_attainment
% lean: CausalSmith.Stat.LdpOptvalueUniformFrontier.causalAttainment_mono_constant

\item
Suppose \(t>1\). We construct the three transcript statistics used in this branch. Define the binary and vector block sizes, privacy attenuation, and scaled-message magnitude by
\[
m_{\mathrm{bin}}=\left\lfloor\frac n2\right\rfloor,
\qquad
m_{\mathrm{vec}}=n-m_{\mathrm{bin}},
\qquad
\delta=\tanh(\varepsilon/2)
=\frac{e^\varepsilon-1}{e^\varepsilon+1},
\qquad
b=\frac2\delta.
\]
Here \(m_{\mathrm{bin}}\ge1\), \(m_{\mathrm{vec}}\ge1\), and \(0<\delta<1\).
For later calibration, the elementary inequalities
\[
e^\varepsilon-1\ge\varepsilon,
\qquad e^\varepsilon\le e<3
\]
also give
\[
\delta
=\frac{e^\varepsilon-1}{e^\varepsilon+1}
\ge\frac{\varepsilon}{4}.
\]

Use a singleton public seed. Each participant's message belongs to the finite alphabet
\(\{-1,1\}\times\{-1,1\}^{2}\). For an original record
\(O=(j,A,Y)\), with \(j\in[2]\) and \(A,Y\in\{0,1\}\), define, for
\(h\in\{-1,1\}\), \(z\in\{-1,1\}^{2}\), and \(1\le i\le n\),
\[
Q_i\bigl(\{(h,z)\}\mid O=(j,A,Y)\bigr)
=
\begin{cases}
\displaystyle
\frac{1+h\delta(2Y-1)}2\,
\mathbf1_{\{z=(-1,-1)\}},
& i\le m_{\mathrm{bin}},\\[6pt]
\displaystyle
\frac14\bigl(1+\delta(2A-1)(2Y-1)z_j\bigr)\,
\mathbf1_{\{h=-1\}},
& i>m_{\mathrm{bin}}.
\end{cases}
\]
These kernels are defined on every original record and are independent of message history. Their masses are nonnegative and sum to one: the first branch sums over the two binary signs, while in the second branch
\[
\sum_{z\in\{-1,1\}^{2}}z_j=0.
\]

For the binary branch, its two possible nonzero atom masses lie between
\((1-\delta)/2\) and \((1+\delta)/2\). The identity
\[
1+\delta=e^\varepsilon(1-\delta)
\]
therefore gives the atomwise privacy inequality for every pair of original records. The vector branch is precisely the \(d=2\) channel in
\cref{obj:lem:private-moment-and-dependence}, which supplies its full-record privacy. The fixed extra message coordinates preserve these inequalities. Summing over atoms gives
\[
Q_i(E\mid o)\le e^\varepsilon Q_i(E\mid o')
\]
for every message event \(E\) and every pair of original records \(o,o'\). Thus
\[
Q\in\mathfrak Q_{\mathrm{NI}}(n,2,\varepsilon).
\]

Write the transcript and its scaled vector coordinates as
\[
\mathbf Z
=\bigl((\mathsf H_i^{\mathrm{cal}},Z_i^{\mathrm{cal}})\bigr)_{i=1}^{n},
\qquad
W_{ij}^{\mathrm{cal}}
=bZ_{m_{\mathrm{bin}}+i,j}^{\mathrm{cal}}
\quad
(1\le i\le m_{\mathrm{vec}},\ j\in[2]).
\]
Define the baseline and contrast estimates by
\[
\widehat B_{\mathrm{cal}}
=\frac12+
\frac1{2\delta m_{\mathrm{bin}}}
\sum_{i=1}^{m_{\mathrm{bin}}}\mathsf H_i^{\mathrm{cal}},
\qquad
\widehat\tau_j
=\frac1{m_{\mathrm{vec}}}
\sum_{i=1}^{m_{\mathrm{vec}}}W_{ij}^{\mathrm{cal}}
\quad(j\in[2]).
\]

Conditional on the records, the transcript atom probability is the product of its participant-specific atom probabilities. Integrating against the product record law from
\cref{obj:ass:iid-people}, and using
\cref{obj:ass:independent-randomness}, gives the transcript identity
\[
\mathbb P\{\mathbf Z=((h_i,z_i))_{i=1}^{n}\}
=
\prod_{i=1}^{n}
\int Q_i(\{(h_i,z_i)\}\mid o)\,P_O(do).
\]
Thus distinct message rows are independent under the sampling decision law.

For \(i\le m_{\mathrm{bin}}\), direct summation over the binary message gives
\[
\mathbb E[\mathsf H_i^{\mathrm{cal}}\mid O_i]
=\delta(2Y_i-1),
\qquad
\mathbb E[\mathsf H_i^{\mathrm{cal}}]
=\delta(2B(P)-1),
\qquad
\mathbb E[(\mathsf H_i^{\mathrm{cal}})^2]=1.
\]
It follows that
\[
\mathbb E[\widehat B_{\mathrm{cal}}]=B(P),
\qquad
\operatorname{Var}\left(
\sum_{i=1}^{m_{\mathrm{bin}}}\mathsf H_i^{\mathrm{cal}}
\right)
=
\sum_{i=1}^{m_{\mathrm{bin}}}
\operatorname{Var}(\mathsf H_i^{\mathrm{cal}})
\le m_{\mathrm{bin}}.
\]
Scaling this centered sum proves
\[
\mathbb E[(\widehat B_{\mathrm{cal}}-B(P))^2]
\le\frac1{4m_{\mathrm{bin}}\delta^2}.
\]

For the vector block, \cref{obj:lem:private-moment-and-dependence} supplies
\[
\mathbb E[W_{ij}^{\mathrm{cal}}]=\tau_j(P),
\qquad
\mathbb E[(W_{ij}^{\mathrm{cal}})^2]=b^2.
\]
The hypotheses of that result hold with block size \(m_{\mathrm{vec}}\), since
\(0<m_{\mathrm{vec}}\le n\), \(d=2\), and \(P\in\mathcal V_2\).
Expanding the squared column sum and using independence of distinct rows yields, for each \(j\in[2]\),
\[
\begin{aligned}
\mathbb E[\widehat\tau_j]
&=\tau_j(P),\\
\mathbb E[\widehat\tau_j^2]
&=\frac{
m_{\mathrm{vec}}b^2+
m_{\mathrm{vec}}(m_{\mathrm{vec}}-1)\tau_j(P)^2
}{m_{\mathrm{vec}}^2}
=\tau_j(P)^2+
\frac{b^2-\tau_j(P)^2}{m_{\mathrm{vec}}}.
\end{aligned}
\]
Subtracting the squared mean therefore gives
\[
\mathbb E[(\widehat\tau_j-\tau_j(P))^2]
=
\frac{b^2-\tau_j(P)^2}{m_{\mathrm{vec}}}
\le\frac{b^2}{m_{\mathrm{vec}}}
=\frac4{m_{\mathrm{vec}}\delta^2}.
\]
All three statistics are Borel functions of the finite transcript alone. These are the required component bounds, simultaneously for every \(P\in\mathcal V_2\).
% lean: CausalSmith.Stat.LdpOptvalueUniformFrontier.twoCellComponents_of_sampling

\item
In the present branch,
\[
1<\sqrt n\,\varepsilon,
\qquad
\min\left\{1,\frac1{n\varepsilon^2}\right\}
=\frac1t,
\qquad
\min\left\{1,\frac1{\sqrt n\,\varepsilon}\right\}
=\frac1{\sqrt n\,\varepsilon}.
\]
Define the unprojected and projected welfare estimates by
\[
\widetilde V_{\mathrm{cal}}
=\widehat B_{\mathrm{cal}}
+\frac{|\widehat\tau_1|+|\widehat\tau_2|}{4},
\qquad
T(\mathbf Z,R,U)
=\max\left\{\frac14,
\min\left\{\frac34,\widetilde V_{\mathrm{cal}}\right\}\right\}.
\]
This \(T\) is a Borel function of \(\mathbf Z\) alone.
By \cref{obj:prop:signed-causal-bridge},
\[
V(P)=B(P)+\frac{|\tau_1(P)|+|\tau_2(P)|}{4}.
\]
The interior means also give \(V(P)\in[1/4,3/4]\), as established in the preceding branch.

Projection toward this interval reduces squared error:
\[
(T-V(P))^2
\le(\widetilde V_{\mathrm{cal}}-V(P))^2.
\]
Indeed, if \(\widetilde V_{\mathrm{cal}}\le1/4\), then
\(T=1/4\le V(P)\); if \(\widetilde V_{\mathrm{cal}}\ge3/4\), then
\(V(P)\le T=3/4\); and inside the interval \(T=\widetilde V_{\mathrm{cal}}\).

The reverse triangle inequality gives
\[
\bigl(|\widehat\tau_j|-|\tau_j(P)|\bigr)^2
\le(\widehat\tau_j-\tau_j(P))^2
\qquad(j\in[2]).
\]
Applying the squared-sum inequality first to the baseline and contrast contributions and then to the two contrast contributions proves the pointwise bound
\[
(T-V(P))^2
\le
2(\widehat B_{\mathrm{cal}}-B(P))^2
+\frac{
(\widehat\tau_1-\tau_1(P))^2+
(\widehat\tau_2-\tau_2(P))^2
}{4}.
\]
Integrating and substituting the component bounds gives
\[
\mathbb E[(T-V(P))^2]
\le
\frac1{2m_{\mathrm{bin}}\delta^2}
+\frac2{m_{\mathrm{vec}}\delta^2}.
\]

The half split satisfies
\[
m_{\mathrm{bin}}\ge\frac n3,
\qquad
m_{\mathrm{vec}}\ge\frac n2.
\]
For the first inequality, even \(n\) gives \(m_{\mathrm{bin}}=n/2\); odd \(n\ge3\) gives
\(m_{\mathrm{bin}}=(n-1)/2\ge n/3\).
The second follows from \(m_{\mathrm{bin}}\le n/2\).
Together with \(\delta\ge\varepsilon/4\), these bounds yield the numerical calibration
\[
\begin{aligned}
\mathbb E[(T-V(P))^2]
&\le
\frac1{2(n/3)\delta^2}
+\frac2{(n/2)\delta^2}\\
&=\frac{88}{n(4\delta)^2}
\le\frac{88}{n\varepsilon^2}
=\frac{88}{t}
\le\frac{1000}{t}.
\end{aligned}
\]
% lean: CausalSmith.Stat.LdpOptvalueUniformFrontier.two_cell_risk_of_component_mse

\item
Using the same protocol and estimator, define the deterministic interval radius and interval by
\[
q_{\mathrm{cal}}
=\sqrt{\frac{10000}{t}}
=\frac{100}{\sqrt n\,\varepsilon}>0,
\]
\[
I(\mathbf Z,R,U)
=
\left[
\max\left\{\frac14,T-q_{\mathrm{cal}}\right\},
\min\left\{\frac34,T+q_{\mathrm{cal}}\right\}
\right].
\]
Both endpoints are Borel functions of the transcript. Since
\(T\in[1/4,3/4]\) and \(q_{\mathrm{cal}}\ge0\),
\[
\frac14
\le\max\left\{\frac14,T-q_{\mathrm{cal}}\right\}
\le T
\le\min\left\{\frac34,T+q_{\mathrm{cal}}\right\}
\le\frac34.
\]
Thus the endpoints are ordered and the interval is connected, closed, and contained in the required domain.

For \(V(P)\in[1/4,3/4]\), the inequality
\(|T-V(P)|\le q_{\mathrm{cal}}\) implies \(V(P)\in I\). Hence
\[
\{V(P)\notin I\}
\subseteq
\{(T-V(P))^2\ge q_{\mathrm{cal}}^2\}.
\]
Markov's inequality and the risk bound give
\[
\begin{aligned}
q_{\mathrm{cal}}^2\,
\mathbb P\{(T-V(P))^2\ge q_{\mathrm{cal}}^2\}
&\le\mathbb E[(T-V(P))^2]\\
&\le\frac{88}{t}
\le\frac{1000}{t}
=\frac{q_{\mathrm{cal}}^2}{10}.
\end{aligned}
\]
Division by \(q_{\mathrm{cal}}^2>0\) and passage to the coverage complement therefore yield
\[
\mathbb P\{V(P)\in I\}\ge1-\frac1{10}=0.90.
\]

Finally, the upper endpoint is at most \(T+q_{\mathrm{cal}}\), and the lower endpoint is at least \(T-q_{\mathrm{cal}}\). Consequently,
\[
\operatorname{length}I
\le2q_{\mathrm{cal}}
=\frac{200}{\sqrt n\,\varepsilon},
\]
and, under the probability decision law,
\[
\mathbb E[\operatorname{length}I]
\le\frac{200}{\sqrt n\,\varepsilon}
\le\frac{1000}{\sqrt n\,\varepsilon}.
\]
This completes the simultaneous attainment bounds in the branch \(t>1\).
% lean: CausalSmith.Stat.LdpOptvalueUniformFrontier.twoCellInterval_coverage_of_risk
% lean: CausalSmith.Stat.LdpOptvalueUniformFrontier.twoCellInterval_expectedLength_le
% lean: CausalSmith.Stat.LdpOptvalueUniformFrontier.two_cell_unsaturated_attainment_of_components

\item
The two branches provide, for every allowed tuple and sampling scheme, a total noninteractive protocol on the original records and transcript-based decisions satisfying
\[
\sup_{P\in\mathcal V_2}\mathbb E[(T-V(P))^2]
\le1000\min\left\{1,\frac1{n\varepsilon^2}\right\},
\]
\[
\inf_{P\in\mathcal V_2}\mathbb P\{V(P)\in I\}\ge0.90,
\qquad
\sup_{P\in\mathcal V_2}\mathbb E[\operatorname{length}I]
\le1000\min\left\{1,\frac1{\sqrt n\,\varepsilon}\right\}.
\]
The definitions in
\cref{obj:def:noninteractive-class,obj:def:sequential-class} give
\[
\mathfrak Q_{\mathrm{NI}}(n,2,\varepsilon)
\subseteq
\mathfrak Q_{\mathrm{SI}}(n,2,\varepsilon).
\]
Thus these same witnesses are admissible for either protocol class. Taking the infima in
\cref{obj:def:risk,obj:def:honest-length} proves, for every
\(\mathcal C\in\{\mathrm{NI},\mathrm{SI}\}\),
\[
\mathfrak R_{\mathcal C}(n,2,\varepsilon)
\le1000\min\left\{1,\frac1{n\varepsilon^2}\right\},
\qquad
\mathfrak H_{\mathcal C}(n,2,\varepsilon)
\le1000\min\left\{1,\frac1{\sqrt n\,\varepsilon}\right\}.
\]
% lean: CausalSmith.Stat.LdpOptvalueUniformFrontier.causalAttainment_minimax_upper

\item
Apply the lower bounds obtained at the start to the fixed sampling scheme and either protocol class. Since \(c\le c_1\) and both rate factors are nonnegative,
\[
c\min\left\{1,\frac1{n\varepsilon^2}\right\}
\le
c_1\min\left\{1,\frac1{n\varepsilon^2}\right\}
\le\mathfrak R_{\mathcal C}(n,2,\varepsilon),
\]
\[
c\min\left\{1,\frac1{\sqrt n\,\varepsilon}\right\}
\le
c_1\min\left\{1,\frac1{\sqrt n\,\varepsilon}\right\}
\le\mathfrak H_{\mathcal C}(n,2,\varepsilon).
\]
Together with the upper bounds and the simultaneous witnesses, these prove all assertions with the chosen universal constants \(c\) and \(C_0\).
% lean: CausalSmith.Stat.LdpOptvalueUniformFrontier.two_cell_calibration
\end{enumerate}
\end{proof}
